% ================================================================
% BLL-VORLAGE: HIER ALLE PERSONEN- UND SCHULDATEN EINTRAGEN
% ================================================================
% Alles in geschweiften Klammern darf geändert werden.
% Bitte die Beispielfelder vor Abgabe vollständig ersetzen!

\newcommand{\Arbeitstitel}{Untersuchung der Schwingungsdauer eines Fadenpendels}
\newcommand{\Untertitel}{Theoretische Grundlagen und Konzept einer experimentellen Prüfung}
\newcommand{\Kurztitel}{Schwingungsdauer eines Fadenpendels}

\newcommand{\Autor}{Vorname Nachname}
\newcommand{\Schule}{Name der Schule}
\newcommand{\Schulort}{Kassel}
\newcommand{\Jahrgang}{Qualifikationsphase Q3/Q4}
\newcommand{\KlasseOderKurs}{Physik (Referenzfach)}
\newcommand{\Abiturjahr}{2027}
\newcommand{\BetreuungSchule}{Name der betreuenden Lehrkraft}
\newcommand{\Abgabedatum}{TT. Monat JJJJ}

% Externe Begleitung ist optional und KEINE schulische Prüfungsfunktion.
% Wenn keine externe Betreuung besteht: \MitSFNfalse setzen.
\newif\ifMitSFN
\MitSFNtrue
\newcommand{\ExterneEinrichtung}{Schülerforschungszentrum Nordhessen (SFN)}
\newcommand{\BetreuungSFN}{Name der fachlichen Ansprechperson}

% Falls die Schule keinen abstrakten Text / keine Verzeichnisse verlangt,
% die betreffenden \input-Zeilen in bll.tex entfernen.
% Logos sind optional: logos/schule.pdf oder logos/schule.png
%                      logos/sfn.pdf oder logos/sfn.png
